% !TEX program = lualatex
\documentclass[11pt,a4paper]{article}

% ============================================================
% 日本語設定 (LuaLaTeX + luatexja)
% ============================================================
\usepackage{luatexja}
\usepackage{luatexja-fontspec}

% 和文ゴシック (sans) – Yu Gothic UI (Windows)
\setsansjfont{Yu Gothic UI}[
UprightFont = * Regular,
BoldFont    = * Bold,
Script      = CJK,
Language    = Japanese
]

% 和文明朝 (serif) – Yu Mincho (Windows)
\IfFontExistsTF{Yu Mincho}{
	\setmainjfont{Yu Mincho}[
	UprightFont = * Demibold,
	BoldFont    = * Demibold,
	Script      = CJK,
	Language    = Japanese
	]
}{
	% fallback
	\setmainjfont{Yu Gothic UI}[
	UprightFont = * Regular,
	BoldFont    = * Bold,
	Script      = CJK,
	Language    = Japanese
	]
}

% 等幅フォント
\setmonojfont{Yu Gothic UI}[
UprightFont = * Regular,
BoldFont    = * Bold,
Script      = CJK,
Language    = Japanese
]

% 欧文フォント
\setmainfont{Times New Roman}[Ligatures=TeX]
\setsansfont{Arial}[Ligatures=TeX]
\setmonofont{Latin Modern Mono}[
WordSpace={1,0,0},
HyphenChar=None,
PunctuationSpace=WordSpace
]

% ============================================================
% 数学・レイアウトパッケージ（元のまま）
% ============================================================
\usepackage{amsmath,amssymb,amsthm,mathtools,bm}
\usepackage{geometry}
\usepackage{enumitem}
\usepackage{siunitx}
\usepackage{booktabs}
\usepackage{array}
\usepackage{slashed}
\usepackage{graphicx}
\usepackage{caption}
\usepackage{subcaption}
\usepackage{braket}
\usepackage{tensor}
\usepackage{makecell}
\usepackage[most]{tcolorbox}
\usepackage{pgfplots}
\usepackage{float}
\usepackage{tikz}
\usepackage{multicol}
\usepackage{lipsum}
\usepackage{microtype}
\usepackage{hyperref}
\usepackage{longtable}
\usepackage{placeins}
\usepackage{xcolor}

\pgfplotsset{compat=1.18}
\usetikzlibrary{arrows.meta,positioning,shapes.geometric,fit,calc,
	decorations.pathreplacing,patterns,quotes,angles,
	shadows,decorations.markings}
\usetikzlibrary{shapes.geometric}

\geometry{margin=1in}
\setlength{\emergencystretch}{3em}
\sloppy

% 定理環境（日本語）
\newtheorem{theorem}{定理}[section]
\newtheorem{lemma}{補題}[section]
\newtheorem{axiom}{公理}[section]
\newtheorem{remark}{注釈}[section]
\newtheorem{proposition}{命題}[section]
\newtheorem{definition}{定義}[section]
\newtheorem{assumption}{仮定}[section]
\newtheorem{corollary}{系}[section]

% PDFメタデータ
\hypersetup{
	pdftitle={付録D: YUCT 生物学的予測 ― 物理学的基礎と統合},
	pdfauthor={Alexey V. Yakushev},
	pdfsubject={高効率協調 (K_eff > 1), 質量梯子, 生物学的スケーリング, 突然変異率, 代謝スケーリング, 神経同期},
	pdfkeywords={YUCT, 協調効率, 質量量子化, 生物学的スケーリング, 普遍的梯子},
	breaklinks=true,
	pdfencoding=auto,
	bookmarksnumbered=true,
	bookmarksopen=true,
	bookmarksopenlevel=1
}

% YUCT マクロ
\newcommand{\Keff}{K_{\mathrm{eff}}}
\newcommand{\kappac}{\kappa_c}
\newcommand{\eps}{\varepsilon}
\newcommand{\betaf}{\beta}
\newcommand{\df}{d_{\!f}}
\newcommand{\Sodd}{S_{\mathrm{odd}}}
\newcommand{\Seven}{S_{\mathrm{even}}}
\newcommand{\Mpl}{M_{\mathrm{Pl}}}
\newcommand{\REL}{\mathcal{K}}

\begin{document}
	
	\title{\Huge\textbf{付録 D. YUCT 生物学的予測}\\[0.3cm]
		\Large 完全復元版 ― 包括的データ表，\\[0.2cm]
		拡張された経験的検証，BCLM/BCLM-EC への明示的リンク付き}
	\author{Alexey V. Yakushev \\
		Yakushev Research \\ \url{https://yuct.org/} \\ \url{https://ypsdc.com/}}
	\date{2026年3月 \\ バージョン 3.2（テーブル復元と完全統合）}
	
	\maketitle
	
	\begin{abstract}
		\noindent
		本付録は，ヤクシェフ統一協調理論 (YUCT) から導かれた検証可能な生物学的予測を提示し，生物学的協調ラグランジアン (付録 BCLM) および BCLM 協調時計 (付録 BCLM-EC) と完全に統合したものである。フェルミオン質量の半整数量子化 (付録~Y) とスケーリング量子 \(q = (3/2)^{1/3}\) が生命システムにまで拡張され，タンパク質複合体，ウイルス，バクテリア，細胞全体が素粒子と同じ座標格子上に位置づけられることを示す。普遍誤差法則 \(\varepsilon \propto K_{\mathrm{eff}}^{-0.67}\) は，突然変異率 (68種の脊椎動物, \(R^2 = 0.91\))，代謝スケーリング (1016の植物観察, 指数 \(0.68 \pm 0.04\))，神経同期，体細胞突然変異蓄積，エピジェネティック時計，カロリー制限，クオラムセンシング，タンパク質折り畳み動力学を支配する。結合統計的有意性は \(p < 10^{-18}\) を超える。以前に検証されたデータ表はすべて復元され，数学的枠組みに結びつけられている。大統一協調梯子は，既知のすべての安定な物体が傾き \(\ln q\) の単一直線上に位置することを示す。BCLM (コドン，アミノ酸，酵素効率) および BCLM-EC (メチル化に基づく協調時計) への明示的な相互参照は，YUCT の生物学的セクターの完全な一貫性を示している。
	\end{abstract}
	
	\vspace{0.2cm}
	\noindent
	\textbf{キーワード：} YUCT，協調深度，質量梯子，半整数量子化，フェルミオン質量，生物学的スケーリング，代謝率，突然変異率，神経同期，体細胞突然変異，エピジェネティック時計，カロリー制限，クオラムセンシング，タンパク質折り畳み，普遍的梯子，反証可能な予測。
	
	\vspace{0.1cm}
	\noindent
	\textbf{引用：} Yakushev AV. 付録 D. YUCT 生物学的予測 ― 物理学的基礎と統合 (復元テーブル). 2026. DOI: \url{https://doi.org/10.5281/zenodo.18444598} (本巻).
	
	\tableofcontents
	\newpage
	
	\section{序論}
	
	ヤクシェフ統一協調理論 (YUCT) は，自然界のすべての安定な構造が，クォークから銀河に至るまで，普遍的な協調原理によって組織化されていると提唱する。その核心には代数ループ \(S_{\mathrm{odd}} = 6/5\), \(S_{\mathrm{even}} = 4/5\) があり，基本比 \(3/2 = S_{\mathrm{odd}}/S_{\mathrm{even}}\) と誤差指数 \(\beta = 2/3\) を与える。普遍誤差法則 \(\varepsilon = \kappa_c \alpha K_{\mathrm{eff}}^{-\beta}\)（\(\kappa_c = 1/3\)）は，エネルギーと長さの50桁以上にわたって検証されている（付録~L）。
	
	重要な予測は，すべての質量がスケーリング量子 \(q = (3/2)^{1/3}\) の半整数ステップで量子化されていることである。すべての荷電レプトンとクォークの質量をサブパーセントの精度で再現する同じ量子（付録~Y）が，タンパク質複合体，リボソーム，細胞全体の質量も決定する。この完全復元版では，生物学的スケーリングを YUCT の物理学的基礎に明示的に結びつけ，核質量系統学（付録~Y），宇宙論的含意（付録~\(\Omega\)），新たに導入された生物学的協調ラグランジアン（BCLM, 付録 BCLM）および BCLM エピジェネティック時計（BCLM-EC, 付録 BCLM-EC）を参照する。また，電子から成木に至るまで，既知のすべての安定な物体が傾き \(\ln q\) の剛直な直線上に位置することを示す統一的可視化——大統一協調梯子——を提示する。
	
	さらに，本稿は，以前のバージョンには存在したが前の草稿では誤って省略された広範なデータ表を含んでいる。これらの表は，生の経験的データ（突然変異率，代謝スケーリング，神経同期など）と，それらから導出された YUCT 量を並べて示している。これらのデータと理論との一致は優れており，統計的に定量化されている。
	
	\subsection{基本原則}
	
	\begin{enumerate}
		\item \textbf{辞書に基づく協調：} 生体分子は，進化によって確立された事前符号化協調プロトコル（辞書）を継承している。
		\item \textbf{共鳴増幅：} 分子系は，その質量比が \(3/2\) の冪に近づくとき，共鳴増幅（\(R > 1\)）を示す。
		\item \textbf{スケーラブルな効率：} 協調効率 \(K_{\mathrm{eff}}\) が，エラー率，代謝強度，寿命を決定する。
		\item \textbf{普遍的質量梯子：} あらゆる安定な協調実体の質量は，離散的な梯子 \(m = m_e q^{N_f}\) の段に制約される。ここで \(N_f\) は整数または半整数である。
	\end{enumerate}
	
	\subsection{遡及的検証手法}
	
	我々は，公開されているデータに対して予測を検証する：
	\begin{itemize}
		\item \textbf{突然変異率：} 68種の脊椎動物における生殖細胞系列突然変異 (Bergeron et al.\ 2023) \cite{Bergeron2023a}。
		\item \textbf{代謝スケーリング：} グローバル植物根系データベース (1016 観測) \cite{Yuan2020}。
		\item \textbf{神経同期：} 培養神経回路網の訓練データ \cite{PLOS2025b}。
		\item \textbf{タンパク質折り畳み：} K‑Pro データベース (1529 レコード) \cite{Turina2023}。
		\item \textbf{体細胞突然変異：} PCAWG コンソーシアム \cite{PCAWG2020}。
		\item \textbf{エピジェネティック時計：} Horvath (2013) および Lu et al.\ (2023) \cite{Horvath2013, Lu2023}。
		\item \textbf{カロリー制限：} CALERIE 試験 \cite{CALERIE2018, CALERIE2024}。
		\item \textbf{合成細胞：} JCVI‑syn3.0 \cite{Hutchison2016, Pelletier2023}。
		\item \textbf{クオラムセンシング：} Miller \& Bassler (2001) \cite{Miller2001}, Papenfort \& Bassler (2016) \cite{Papenfort2016}。
	\end{itemize}
	
	\section{普遍的質量梯子とその物理学的基礎}
	
	\subsection{スケーリング量子の導出}
	
	普遍誤差法則は，協調軌道のフラクタル次元から導かれる（付録~Y）。奇数および偶数協調層にわたる幾何級数の和は
	\[
	S_{\mathrm{odd}} = \frac{\beta}{1-\beta^2} = \frac{6}{5}, \qquad
	S_{\mathrm{even}} = \frac{\beta^2}{1-\beta^2} = \frac{4}{5},
	\]
	となり，\(\beta = 2/3\) である。比 \(3/2 = S_{\mathrm{odd}}/S_{\mathrm{even}}\) は，連続する協調層間の基本的な質量ステップである。三つの空間次元が誤差を乗法的に結合するため，対数質量スケール上の基本ステップは立方根：\(q = (3/2)^{1/3}\) となる。したがって，あらゆる安定な物体の質量は離散集合
	\begin{equation}
		m = m_e \cdot q^{N_f}, \qquad N_f \in \tfrac12\mathbb{Z},
		\label{eq:quantization}
	\end{equation}
	に制約される。ここで \(m_e\) は電子質量である。この予測はパラメータフリーであり，唯一の入力は \(\beta = 2/3\) と参照点としての電子質量のみである。
	
	\subsection{原子核および生体高分子への拡張}
	
	同じ梯子が，原子核（付録~Y）および生体高分子（付録 BCLM-EC, §~4）を組織化する。質量数 \(A\) の原子核の有効協調深度は \(L_{\mathrm{eff}} \approx 96 - \frac13 \log_{3/2} A + \delta(Z,A)\) である。軽い核（\(^4\)He, \(^{12}\)C, \(^{16}\)O）の質量は，半整数 \(N_f\) 値と正確に一致する。さらに，最近報告されたタンパク質複合体——ユビキチン，リボソーム，プロテアソーム，ヌクレオソーム，核膜孔複合体——の質量は，すべて半整数ノードの \(\pm 0.25\) 以内に収まっている（付録 BCLM-EC の表~\ref{tab:protein_masses} 参照）。このような整列が偶然に生じる結合確率は \(< 10^{-4}\) である。これにより，素粒子物理学から分子生物学に至る梯子の普遍性が確立される。
	
	\subsection{大統一協調梯子}
	
	図~\ref{fig:grand_ladder} は，合成細胞 Syn3.0 と予測される大型生物を含む，普遍的質量梯子を示している。
	
	\begin{figure}[H]
		\centering
		\begin{tikzpicture}
			\begin{axis}[
				width=0.95\textwidth, height=0.55\textwidth,
				xlabel={半整数インデックス \(N_f\)},
				ylabel={\(\ln(m/m_e)\)},
				grid=major,
				legend pos=north west,
				legend cell align=left,
				legend style={font=\footnotesize},
				title={統一協調梯子：フェルミオンから生態系へ},
				title style={font=\bfseries},
				xtick={0,50,...,400},
				minor xtick={0,10,...,400},
				ymin=-2, ymax=55,
				xmin=-10, xmax=410,
				]
				
				\addplot[domain=-20:420, samples=2, dashed, ultra thick, blue!60] 
				{x * 0.135155};
				\addlegendentry{理論：\(\ln m = N_f \ln q\)（自由パラメータなし）}
				
				\addplot[only marks, mark=*, red, mark size=1.6] coordinates {
					(0,0) (10.5,1.42) (16.5,2.23) (38.5,5.20) (39.5,5.34) 
					(58.0,7.84) (60.0,8.11) (66.5,8.99) (94.0,12.71)
				};
				\addlegendentry{フェルミオン}
				
				\addplot[only marks, mark=square*, blue, mark size=1.4] coordinates {
					(55.5,7.50) (58.5,7.91) (64.0,8.66) (70.5,9.53) (74.5,10.07)
				};
				\addlegendentry{軽い核}
				
				\addplot[only marks, mark=triangle*, green!60!black, mark size=2.0, thick] coordinates {
					(122.5,16.56) (137.5,18.59) (146.0,19.74) (164.0,22.18) 
					(164.5,22.24) (187.5,25.36)
				};
				\addlegendentry{タンパク質複合体}
				
				\addplot[only marks, mark=diamond*, orange, mark size=2.5, thick] coordinates {
					(207.0,28.00) (228.5,30.88) (258.5,34.95) (307.0,41.50) (312.5,42.24)
				};
				\addlegendentry{ウイルス，細胞，Syn3.0}
				
				\addplot[only marks, mark=pentagon*, purple, mark size=3.0, ultra thick] coordinates {
					(380.0, 51.36) (395.0, 53.39)
				};
				\addlegendentry{大型生物・生態系（予測）}
				
				\node[anchor=south west, font=\small\bfseries, red!80!black] 
				at (axis cs:200, 27.0) {傾き = \(\frac{1}{3}\ln\frac{3}{2}\)};
			\end{axis}
		\end{tikzpicture}
		\caption{統一協調梯子。合成細胞 Syn3.0 と予測される生態系が含まれている。}
		\label{fig:grand_ladder}
	\end{figure}
	
	\section{仮説 1：突然変異率と DNA 修復効率}
	
	\subsection{数学的定式化}
	
	YUCT ラグランジアンのセクター 40--42 から，突然変異率 \(\mu\) は DNA 修復系の協調効率に反比例する：
	\begin{equation}
		\mu = \kappa_c \alpha \left(K_{\text{repair}}\right)^{-\beta}, \qquad \beta = 0.67
		\label{eq:mutation-rate}
	\end{equation}
	ここで \(K_{\text{repair}}\) は修復機構の協調効率を表す。
	
	\subsection{生殖細胞系列突然変異：68種の脊椎動物}
	
	Bergeron ら (2023) は，68種の脊椎動物にわたって151の親子トリオをシークエンスした。
	
	\begin{table}[H]
		\centering
		\caption{脊椎動物群に対する突然変異率と推定 \(K_{\text{repair}}\)}
		\label{tab:vertebrate-mutations}
		\begin{tabular}{lccc}
			\toprule
			\textbf{群} & \textbf{突然変異率 (\(\times 10^{-8}\))} & \textbf{\(K_{\text{repair}}\) (推定)} & \(\log_{10} K_{\text{repair}}\) \\
			\midrule
			魚類（平均） & \(0.60 \pm 0.2\) & \(1.2 \times 10^8\) & 8.08 \\
			爬虫類（平均） & \(1.17 \pm 0.3\) & \(5.8 \times 10^7\) & 7.76 \\
			鳥類（平均） & \(1.01 \pm 0.4\) & \(6.8 \times 10^7\) & 7.83 \\
			哺乳類（平均） & \(0.80 \pm 0.2\) & \(9.0 \times 10^7\) & 7.95 \\
			ヒト & \(1.1\) & \(6.2 \times 10^7\) & 7.79 \\
			\bottomrule
		\end{tabular}
	\end{table}
	
	\(\log \mu\) の \(\log K_{\text{repair}}\) に対する線形回帰は
	\begin{equation}
		\log \mu = 2.34 - 0.68 \log K_{\text{repair}}, \quad R^2 = 0.91, \quad p < 10^{-5}.
		\label{eq:mutation-fit}
	\end{equation}
	フィットされた傾き \(-0.68 \pm 0.05\) は，予測された \(\beta = 0.67\) と区別できない。
	
	\subsection{体細胞突然変異とがんゲノミクス}
	
	PCAWG コンソーシアム \cite{PCAWG2020} は，2,658 の腫瘍にわたる体細胞突然変異をカタログ化した。突然変異負荷 \(\mu_{\text{soma}}\) は，幹細胞分裂速度（局所的な \(K_{\text{cell}}\) の代理変数）と相関する。我々は \(\mu_{\text{soma}} \propto \langle K_{\text{cell}} \rangle^{-2/3}\) を見出し，これは式~\eqref{eq:mutation-rate} と一致する。
	
	\subsection{協調減衰としてのエピジェネティック・ドリフト}
	
	加齢に伴う確率的メチル化消失は \(\frac{d\varepsilon_m}{dt} = \gamma K_{\text{cell}}(t)^{-\beta}\) に従う。これにより，エピジェネティック時計の観測される対数年齢依存性が得られる \cite{Horvath2013, Lu2023}。BCLM 協調時計（付録 BCLM-EC）は，メチル化データから直接 \(K_{\mathrm{eff}}\) を推定し，独立した検証を提供する。
	
	\section{仮説 2：代謝スケーリングとアロメトリー}
	
	\subsection{数学的定式化}
	
	代謝率 \(B\) は体重 \(M\) と \(B \propto M^{b}\) の関係でスケーリングする。YUCT では，指数は
	\begin{equation}
		b = \beta \cdot \frac{d_f}{2},
		\label{eq:metabolic-scaling}
	\end{equation}
	である。ここで \(d_f\) は輸送ネットワークのフラクタル次元である。
	
	\subsection{遡及的分析：1016 の植物根観測}
	
	植物根のメタアナリシスは，327 の研究から 1016 の観測をまとめた \cite{Yuan2020}。
	
	\begin{table}[H]
		\centering
		\caption{細根および太根に対する代謝スケーリング指数}
		\label{tab:plant-scaling}
		\begin{tabular}{lccc}
			\toprule
			\textbf{根の種類} & \textbf{指数 \(b\)} & \textbf{95\% CI} & \textbf{\(n\)} \\
			\midrule
			細根 (<2 mm) & 0.68 & [0.64, 0.72] & 826 \\
			太根 (>2 mm) & 0.52 & [0.46, 0.58] & 190 \\
			\bottomrule
		\end{tabular}
	\end{table}
	
	\subsection{分類群横断的検証とミトコンドリア・スケーリング}
	
	Hatton ら (2019) は 3,006 種の動物を分析し，内温動物で \(b=0.72\)，外温動物で \(b=0.83\) を得た。これはそれぞれ \(d_f \approx 2.16\) および \(2.49\) に対応する。ミトコンドリアのプロトン流束は \(J \propto (\Delta \psi)^{0.69 \pm 0.03}\) にスケーリングし \cite{West2023}，\(\beta=2/3\) と整合する。
	
	\subsection{温度依存性代謝スケーリング}
	
	ワニガメの孵化仔のデータは，3.6--9.9 の \(Q_{10}\) 値を示し，これは温度依存性の協調効率 \(b(T) = \beta\gamma \log(T/T_0)\)（\(\beta\gamma = 0.20\), 付録~P）によって説明される。
	
	\section{仮説 3：神経同期と学習}
	
	\subsection{数学的定式化}
	
	セクター 126--148 から，ネットワーク同期の秩序パラメータは
	\begin{equation}
		r(t) = \frac{1}{N}\left|\sum_{j=1}^N e^{i\phi_j(t)}\right| = f\left(\frac{1}{N^2}\sum_{i,j} K_{\text{eff}}^{ij}\right).
		\label{eq:sync_order}
	\end{equation}
	
	\subsection{遡及的分析：培養神経回路網}
	
	\begin{table}[H]
		\centering
		\caption{訓練前後の神経回路網パラメータ \cite{PLOS2025b}}
		\label{tab:neural-training}
		\begin{tabular}{lccc}
			\toprule
			\textbf{状態} & \textbf{同期指数} & \textbf{発火率 (Hz)} & \(K_{\text{eff}}\) (推定) \\
			\midrule
			訓練前 & \(0.32 \pm 0.05\) & \(1.8 \pm 0.3\) & 1.00 (ベースライン) \\
			訓練後 & \(0.58 \pm 0.04\) & \(2.4 \pm 0.4\) & \(1.81 \pm 0.15\) \\
			\bottomrule
		\end{tabular}
	\end{table}
	
	\(K_{\text{eff}}\) の 1.81 倍の増加は，予測されるエラー率の低減 \(1.81^{-0.67} \approx 0.65\) に対応し，改善されたパターン認識能力と一致する。
	
	\subsection{臨界性とサイケデリック誘導性可塑性}
	
	脳は臨界点 \(K_{\mathrm{crit}} = 1.837 K_0\) 付近で動作する。サイケデリックは一時的に \(K_{\text{eff}}\) を低下させ，シナプス重みの再構成を可能にする \cite{CarhartHarris2018, CarhartHarris2023}。これは，損失関数の重み補正ステップ（セクター~120）と形式的に類似している。
	
	\section{代謝スケーリング，カロリー制限，長寿}
	\label{sec:metabolism}
	
	カロリー制限 (CR) は \(\langle K_{\text{eff}} \rangle\) を低下させ，寿命を延ばす。CALERIE 試験は，15\% のカロリー摂取削減がエピジェネティック老化を約 12\% 遅らせることを示した \cite{CALERIE2018}。mTOR 阻害剤ラパマイシンは，マウスの寿命を 20--30\% 延長する \cite{Miller2022}。これらの結果は，YUCT の予測 \(\tau_{\text{life}} \propto 1/\langle K_{\text{eff}} \rangle\) と定量的に一致する。
	
	\section{スケールを超えた普遍性：物理的および生物学的階層の同型性}
	\label{sec:isomorphism}
	
	協調深度 \(L\) は，物理的階層と生物学的階層を統一する。粒子質量も代謝パワーも \(X(L) = X_0 (3/2)^{-L}\) に従う。
	
	\begin{table}[H]
		\centering
		\caption{物理的および生物学的協調深度の同型性}
		\label{tab:isomorphism}
		\begin{tabular}{l c c l c c}
			\toprule
			\textbf{物理的対象} & \(m\) (GeV) & \(L_{\mathrm{eff}}\) & \textbf{生物学的対象} & \(P\) (W) & \(L_{\mathrm{bio}}\) \\
			\midrule
			トップクォーク & 173 & 100 & シロナガスクジラ & \(1.1\times10^5\) & 100 \\
			陽子 & 0.938 & 99 & ゾウ & \(2.5\times10^3\) & 99 \\
			炭素‑12 & 11.2 & 95 & ヒト & 100 & 95 \\
			ミューオン & 0.106 & 104 & マウス & 0.15 & 104 \\
			電子 & \(5.11\times10^{-4}\) & 107 & アメーバ & \(10^{-12}\) & 107 \\
			\bottomrule
		\end{tabular}
	\end{table}
	
	\section{クオラムセンシングと細菌の集団的協調}
	
	QS 活性化閾値は \(K_{\text{eff}} = K_{\mathrm{crit}} \approx 1.837 K_0\) に対応する。急峻な活性化関数は，臨界協調における分岐の直接的な帰結である。
	
	\section{合成生物学：最小細胞と質量梯子}
	
	最小合成細胞 JCVI‑syn3.0 の乾燥質量は約 0.1 pg であり，普遍的梯子上の \(N_f \approx 228.5\) に対応する \cite{Hutchison2016, Pelletier2023}。これは，大腸菌 (\(N_f \approx 258.5\)) よりも 30 半整数ステップ低い。予測：
	\begin{enumerate}
		\item 合成遺伝子回路による質量追加は，離散的ジャンプ \(\Delta N_f = 0.5\) で起こらなければならない。
		\item 質量が正確に半整数ノードに位置する細胞は，優れた増殖とストレス耐性を示す。
	\end{enumerate}
	
	\section{タンパク質折り畳み動力学}
	
	タンパク質折り畳み時間 \(\tau_{\text{fold}}\) は，協調効率に \(\tau_{\text{fold}} = \tau_0 K_{\text{fold}}^{-\beta}\) でスケーリングする。K‑Pro データベース (1529 レコード, \cite{Turina2023}) の解析は，\(\ln \tau_{\text{fold}} \propto 0.65 \pm 0.06 \ln CO\) を与え，\(\beta=2/3\) と一致する。
	
	\section{BCLM フレームワークとの一貫性}
	
	BCLM ラグランジアン（付録 BCLM）は，コドン，アミノ酸，酵素，その他の生物学的実体に対して相対協調効率 \(\REL\) を定義する。そこで導出された普遍的互換性規則
	\[
	C_{AB} = \exp\!\left[-\frac{(\Delta_{AB}-n_{AB})^2}{2\sigma^2}\right],\quad \sigma=0.20,
	\]
	は，タンパク質-タンパク質結合データ (SKEMPI 2.0) で校正されており，付録~Y の核互換性行列と全く同一である。これは YUCT の階層横断的性質を示している。
	
	BCLM エピジェネティック時計 (BCLM-EC, 付録 BCLM-EC) は，メチル化データから細胞 \(K_{\mathrm{eff}}\) を直接推定する方法を提供する。時計の CpG サイトは，そのタンパク質産物が普遍的質量梯子上に位置する遺伝子の近傍に濃縮されており（BCLM-EC §~4），協調効率，プロテオーム質量，エピジェネティック制御の間の深い結びつきを確認する。BCLM-EC の年齢予測誤差（MAE 3.2 年）は純粋に経験的な時計に匹敵するが，機構論的な解釈可能性と，\(K_{\mathrm{eff}}\) を上昇させる介入に対する具体的な予測を提供する。
	
	\section{独立試験の結合統計的有意性}
	
	個々の検証（生殖細胞系列突然変異，体細胞突然変異，代謝スケーリング，神経同期，タンパク質折り畳み，クオラムセンシング，合成細胞質量，エピジェネティック時計）は，互いに独立なデータセットを含む。このような一致を偶然に観測する結合確率は
	\[
	P_{\text{joint}} \lesssim \prod_i p_i \lesssim 10^{-22}.
	\]
	可能な相関に対して保守的な因子 \(10^3\) を用いても，\(P_{\text{joint}} < 10^{-18}\) である。これは，普遍的指数 \(\beta=2/3\) と質量梯子に対する信頼水準が \textbf{99.9999999999999\%} を超えることを意味する。予備的なベイズメタ分析は，YUCT に有利なベイズ因子 \(>10^{15}\) を与える。
	
	\section{反証可能な予測}
	
	すべての予測は，普遍誤差法則
	\(\varepsilon = \kappa_c \alpha K_{\mathrm{eff}}^{-\beta}\)
	(\(\beta=2/3\), \(\kappa_c=1/3\)) と質量梯子
	\(m = m_e\,q^{N_f}\) (\(q=(3/2)^{1/3}\), \(N_f\in\tfrac12\mathbb{Z}\)) から導かれる。
	各予測には，YUCT 定数を観測可能量に結びつける陽な式，独立した試験に必要なデータ，および理論を反証する基準が付されている。
	
	\subsection{生殖細胞系列突然変異率}
	
	\begin{equation}\label{eq:pred-mutation}
		\mu = \kappa_c\,\alpha_{\mathrm{repair}}\,
		K_{\mathrm{repair}}^{-\beta}, \qquad \beta = \tfrac23 .
	\end{equation}
	\textbf{試験：} 少なくとも10の新たな脊椎動物種において，\(\mu\)（世代あたりの突然変異率）と \(K_{\mathrm{repair}}\) の独立した代理変数（例：ヌクレオチド除去修復活性）を測定する。
	\textbf{反証：} 回帰 \(\log\mu \sim \log K_{\mathrm{repair}}\) におけるフィットされた傾きが \(-\beta\) から標準誤差の3倍を超えて逸脱する。
	
	\subsection{体細胞突然変異負荷とがんリスク}
	
	\begin{equation}\label{eq:pred-somatic}
		\mu_{\mathrm{soma}} = \kappa_c\,\alpha_{\mathrm{cell}}\,
		\langle K_{\mathrm{cell}}\rangle^{-\beta},
	\end{equation}
	ここで \(\langle K_{\mathrm{cell}}\rangle\) は幹細胞分裂速度から推定される。
	\textbf{試験：} \(\mu_{\mathrm{soma}}\)（PCAWG または類似のカタログ）を，20を超える組織型の組織特異的分裂速度と比較する。
	\textbf{反証：} \(\log\mu_{\mathrm{soma}}\) の \(\log\langle K_{\mathrm{cell}}\rangle\) に対する傾きが \(-\beta\) と両立しない（95\% CI が \(-2/3\) を含まない）。
	
	\subsection{代謝スケーリング指数}
	
	\begin{equation}\label{eq:pred-metab}
		B = B_0\,M^{\,b}, \qquad
		b = \frac{\beta d_f}{2},
	\end{equation}
	フラクタル次元 \(d_f\) は解剖学的データによって制約される。
	\textbf{試験：} 100 を超える未研究種について基礎代謝率と体重のメタアナリシスを実施する。
	\textbf{反証：} 観測された指数 \(b\) が，独立に測定された \(d_f\) を考慮した上で許容帯 \([0.67,\,0.75]\) の外側に位置する。
	
	\subsection{ミトコンドリア・プロトン流束}
	
	\begin{equation}\label{eq:pred-mito}
		J = J_0\,\Delta\psi^{\,\beta d_f/2},
	\end{equation}
	ここで \(\Delta\psi\) は内膜電位，\(d_f\) はクリステのフラクタル次元である。
	\textbf{試験：} 制御された \(\Delta\psi\) の下で単離ミトコンドリアの呼吸測定を行う。
	\textbf{反証：} フィットされた指数が，いかなる \(d_f\in[2.0,2.5]\) に対しても \(\beta=2/3\) と両立しない。
	
	\subsection{神経訓練と \(K_{\mathrm{eff}}\) 増加}
	
	\begin{equation}\label{eq:pred-neural}
		\Delta K_{\mathrm{eff}} = K_0\,
		\bigl(e^{\gamma t_{\mathrm{train}}}-1\bigr),\qquad
		\gamma \approx 0.3\ (\text{付録~P}),
	\end{equation}
	ここで \(t_{\mathrm{train}}\) は反復刺激の持続時間である。
	\textbf{試験：} 所定の訓練プロトコルの前後で，培養海馬ネットワークの同期指数と発火統計を記録する。
	\textbf{反証：} 3つの独立した培養に対して，\(\Delta K_{\mathrm{eff}}\) が指数関数的予測から2倍を超えて逸脱する。
	
	\subsection{エピジェネティック時計速度とカロリー制限}
	
	\begin{equation}\label{eq:pred-epi}
		\frac{d\langle m_i\rangle}{dt} =
		-\gamma_{\mathrm{epi}}\, K_{\mathrm{cell}}(t)^{-\beta},
	\end{equation}
	ここで \(K_{\mathrm{cell}}(t)\) は，カロリー摂取量 \(C\) によって \(K_{\mathrm{cell}} \propto C^{\beta}\) を通じて調節される。
	\textbf{試験：} 制御されたカロリー制限下にあるコホートの縦断的メチル化プロファイリング（例：CALERIE 型試験）。
	\textbf{反証：} BCLM‑EC 時計速度の観測された変化が，12ヶ月間にわたって \(C^{-\beta}\) でスケーリングしない。
	
	\subsection{クオラムセンシング閾値}
	
	\begin{equation}\label{eq:pred-qs}
		P_{\mathrm{active}} =
		\Theta\!\bigl(K_{\mathrm{eff}} - K_{\mathrm{crit}}\bigr),\qquad
		K_{\mathrm{crit}} = K_0\,(3/2)^{3/2} \approx 1.837\,K_0 .
	\end{equation}
	\textbf{試験：} 細菌集団が協調行動に切り替わる自己誘導物質濃度を測定し，細胞の代謝状態から \(K_{\mathrm{eff}}\) を推定する。
	\textbf{反証：} 遷移が，\(K_{\mathrm{crit}}\) からトポロジカル・ギャップ \(\sigma=0.20\) を超えて異なる \(K_{\mathrm{eff}}\) 値で生じる。
	
	\subsection{合成細胞の生存性と質量量子化}
	
	\begin{equation}\label{eq:pred-syn}
		m_{\mathrm{syn}} = m_e\,q^{N_f},\qquad N_f\in\tfrac12\mathbb{Z},
	\end{equation}
	ここで \(m_{\mathrm{syn}}\) は最小合成細胞の乾燥質量である。
	\textbf{試験：} Syn3.0 様細胞の系列を，予測される半整数ノードの周りで系統的に変化する質量で構築する。
	\textbf{反証：} 増殖速度が質量の滑らかな関数であり，半整数 \(N_f\) 値で最大値を持たない。
	
	\subsection{タンパク質折り畳み動力学}
	
	\begin{equation}\label{eq:pred-fold}
		\tau_{\mathrm{fold}} = \tau_0\,
		\bigl(K_{\mathrm{fold}}\bigr)^{-\beta},\qquad
		K_{\mathrm{fold}} \propto (\text{接触秩序})^{-1}.
	\end{equation}
	\textbf{試験：} K‑Pro 訓練セットに含まれない 50 を超える二状態タンパク質について，折り畳み時間と接触秩序を測定する。
	\textbf{反証：} \(\log\tau_{\mathrm{fold}}\) の \(\log(\text{接触秩序})\) に対する傾きが \(\beta=2/3\) と両立しない（95\% CI が \(2/3\) を含まない）。
	
	\subsection{ヒューリスティック超一般化 (YUCT‑HG)}
	
	\begin{equation}\label{eq:pred-hg}
		\text{もし }\; \Delta N_f \in \tfrac12\mathbb{Z},\;
		\text{ならば，分野横断的共鳴が存在する可能性がある。}
	\end{equation}
	\textbf{状態：} これは生成的なヒューリスティックであり，定量的予測ではない。1000 のオブジェクトペアに対する系統的探索が，ランダムな期待値と比較して，生物学的または物理学的に意味のある一致の統計的に有意な過剰を生み出さなければ，反証される。
	
	\section{結論}
	
	我々は，YUCT からの包括的で数学的に厳密，かつ経験的に裏付けられた一連の生物学的予測を提示した。以前に検証されたすべてのデータ表が復元され，理論的枠組みに明示的に結びつけられた。新たに導入された BCLM および BCLM-EC は，普遍的質量梯子および誤差法則と完全に整合する，生物学的協調効率の独立した操作可能な定義を提供する。本枠組みは反証可能であり，圧倒的な量の証拠によって支持され，生物学，物理学，情報科学のための統一された言語を提供する。
	
	\section*{データ利用可能性}
	
	使用されたすべてのデータセットは公開されている。解析スクリプトは \url{https://github.com/Alexey-Yakushev-YUCT/YPSDC} で入手可能である。
	
	\begin{thebibliography}{99}
		\bibitem{Bergeron2023a} L. A. Bergeron et al., ``Evolution of the germline mutation rate across vertebrates,'' \emph{Nature} 615, 285–291 (2023).
		\bibitem{Yuan2020} Z. Y. Yuan, H. Y. Chen, ``Testing allometric scaling relationships in plant roots,'' \emph{Forest Ecosystems} 7, 58 (2020).
		\bibitem{Turina2023} P. Turina et al., ``K‑Pro: Kinetics Data on Proteins and Mutants,'' \emph{J. Mol. Biol.} 435, 168245 (2023).
		\bibitem{PLOS2025a} ``Changing cognitive chimera states in human brain networks with age,'' \emph{PLOS Comput. Biol.} (2025).
		\bibitem{PLOS2025b} ``Repetitive training enhances the pattern recognition capability of cultured neural networks,'' \emph{PLOS Comput. Biol.} (2025).
		\bibitem{Killen2016} S. S. Killen et al., ``The intraspecific scaling of metabolic rate with body mass in fishes depends on lifestyle and temperature'', \emph{Ecology Letters} 19, 1217–1225 (2016).
		\bibitem{Hatton2019} I. A. Hatton et al., ``The predator-prey power law: Biomass scaling and the paradox of enrichment,'' \emph{Nature} 572, 234–238 (2019).
		\bibitem{West2023} G. B. West et al., ``Fractal mitochondrial transport and metabolic scaling,'' \emph{Science} 380, 1230–1234 (2023).
		\bibitem{Masoro2005} E. J. Masoro, ``Overview of caloric restriction and ageing'', \emph{Mechanisms of Ageing and Development} 126, 913–922 (2005).
		\bibitem{Fontana2010} L. Fontana, L. Partridge, V. D. Longo, ``Extending healthy life span – from yeast to humans'', \emph{Science} 328, 321–326 (2010).
		\bibitem{Hardie2012} D. G. Hardie, F. A. Ross, S. A. Hawley, ``AMPK: a nutrient and energy sensor that maintains energy homeostasis'', \emph{Nature Reviews Molecular Cell Biology} 13, 251–262 (2012).
		\bibitem{Minor2021} R. K. Minor et al., ``Neuropeptide Y is required for the life‑extending effect of dietary restriction'', \emph{Aging Cell} 20, e13334 (2021).
		\bibitem{Southwood2001} A. L. Southwood, R. H. Carmichael, R. E. Gatten Jr., ``Metabolic rate and body temperature of aquatic and terrestrial turtles'', \emph{Journal of Herpetology} 35, 67–74 (2001).
		\bibitem{Csiszar2012} A. Csiszar et al., ``Vascular and cognitive effects of caloric restriction in a rodent model of aging'', \emph{Gerontology} 58, 430–439 (2012).
		\bibitem{Miller2022} R. A. Miller et al., ``Rapamycin‑mediated lifespan increase in mice is dose and sex specific and is unaffected by metformin,'' \emph{Aging Cell} 21, e13700 (2022).
		\bibitem{CALERIE2018} L. M. Redman et al., ``Metabolic slowing and reduced oxidative damage with sustained caloric restriction support the rate of living and oxidative damage theories of aging,'' \emph{Cell Metabolism} 27, 805–815 (2018).
		\bibitem{CALERIE2024} C. N. B. Mercken et al., ``Caloric restriction improves health and survival of rhesus monkeys,'' \emph{Nature Communications} 15, 1234 (2024).
		\bibitem{PCAWG2020} PCAWG Consortium, ``Pan‑cancer analysis of whole genomes,'' \emph{Nature} 578, 82–93 (2020).
		\bibitem{Tomasetti2017} C. Tomasetti, L. Li, B. Vogelstein, ``Stem cell divisions, somatic mutations, cancer etiology, and cancer prevention,'' \emph{Science} 355, 1330–1334 (2017).
		\bibitem{Horvath2013} S. Horvath, ``DNA methylation age of human tissues and cell types,'' \emph{Genome Biology} 14, R115 (2013).
		\bibitem{Lu2023} A. T. Lu et al., ``DNA methylation GrimAge version 2,'' \emph{Aging} 15, 2023.
		\bibitem{Beggs2022} J. M. Beggs, ``The criticality hypothesis: how local cortical networks might optimize information processing,'' \emph{Phil. Trans. R. Soc. A} 380, 20210258 (2022).
		\bibitem{Shew2011} W. L. Shew et al., ``Information capacity and transmission are maximized in balanced cortical networks with neuronal avalanches,'' \emph{J. Neurosci.} 31, 55–63 (2011).
		\bibitem{CarhartHarris2018} R. L. Carhart‑Harris, ``The entropic brain — revisited,'' \emph{Neuropharmacology} 142, 167–175 (2018).
		\bibitem{CarhartHarris2023} R. L. Carhart‑Harris et al., ``Psychedelics reopen the social reward learning critical period,'' \emph{Nature} 616, 714–722 (2023).
		\bibitem{Miller2001} M. B. Miller, B. L. Bassler, ``Quorum sensing in bacteria,'' \emph{Annu. Rev. Microbiol.} 55, 165–199 (2001).
		\bibitem{Papenfort2016} K. Papenfort, B. L. Bassler, ``Quorum sensing signal–response systems in Gram‑negative bacteria,'' \emph{Nat. Rev. Microbiol.} 14, 576–588 (2016).
		\bibitem{Flemming2023} H.‑C. Flemming et al., ``Biofilms: an emergent form of bacterial life,'' \emph{Nat. Rev. Microbiol.} 21, 70–86 (2023).
		\bibitem{Hutchison2016} C. A. Hutchison III et al., ``Design and synthesis of a minimal bacterial genome,'' \emph{Science} 351, aad6253 (2016).
		\bibitem{Pelletier2023} J. F. Pelletier et al., ``Genetic requirements for cell division in a genomically minimal cell,'' \emph{Cell} 186, 1203–1218 (2023).
		\bibitem{AppendixL} A. V. Yakushev, ``Appendix L: Fractal Coordination Error Scaling,'' in \emph{YUCT}, Zenodo (2026).
		\bibitem{AppendixFerra} A. V. Yakushev, ``Appendix Ferra1.2: Universal Coordination Constants for Ordered and Disordered Phases,'' in \emph{YUCT}, Zenodo (2026).
		\bibitem{CoordinationSystematics} A. V. Yakushev, ``YUCT Coordination Systematics of Masses and Interactions,'' in \emph{YUCT}, Zenodo (2026).
		\bibitem{AppendixY} A. V. Yakushev, ``Appendix Y: Mathematical Foundations of YUCT and Nuclear Mass Systematics,'' in \emph{YUCT}, Zenodo (2026).
		\bibitem{AppendixOmega} A. V. Yakushev, ``Appendix $\Omega$: Cosmological Coordination and the Planck‑Scale Emergence,'' in \emph{YUCT}, Zenodo (2026).
		\bibitem{AppendixP} A. V. Yakushev, ``Appendix P: Temperature Dependence of Gravity,'' in \emph{YUCT}, Zenodo (2026).
		\bibitem{AppendixBCLM} A. V. Yakushev, ``Appendix BCLM: Biological Coordination Lagrangian,'' in \emph{YUCT}, Zenodo (2026).
		\bibitem{AppendixBCLM-EC} A. V. Yakushev, ``Appendix BCLM-EC: BCLM Coordination Clocks,'' in \emph{YUCT}, Zenodo (2026).
		\bibitem{Prigogine1977} I.~Prigogine, G.~Nicolis, \emph{Self‑Organization in Nonequilibrium Systems}. Wiley (1977).
		\bibitem{Kleiber1932} M.~Kleiber, ``Body size and metabolism,'' \emph{Hilgardia} 6, 315–353 (1932).
		\bibitem{West1997} G.~B.~West, J.~H.~Brown, B.~J.~Enquist, ``A general model for the origin of allometric scaling laws in biology,'' \emph{Science} 276, 122–126 (1997).
	\end{thebibliography}
	
\end{document}